\chapter{Field Experience}\label{ch:field-experience}

\emph{What must a candidate show from a training programme of their own, and how is it judged?}
The candidate must show experience of conducting a training programme. Field experience carries 10 marks, part of the practical (140 marks in all). The Board gives the marks and the documents to show. It does not give the detailed content, so ask the institution what it assesses.

\section*{What the candidate must show}

\begin{itemize}
\item A report of a training programme the candidate conducted, with photographs.
\item A feedback form filled by a trainee (Annexure A).
\item The experience details in the Board's prescribed format (Annexure B).
\end{itemize}

Field experience is part of the practical assessment, which lasts at most one hour per candidate.

\section*{Feedback form (Annexure A)}

The form asks the trainee for:
\begin{itemize}
\item name and contact details;
\item period of training and number of days;
\item mode of training: home based, park, society, corporate, other;
\item reason for taking up yoga;
\item feedback about the trainer;
\item a mark out of 10.
\end{itemize}

The form is signed and dated.

\section*{Experience format (Annexure B)}

The format lists:
\begin{itemize}
\item the persons trained, year by year, with the number and the mode of training;
\item working experience with an organisation: name, post, full time, part time or voluntary, dates, hours, role;
\item yoga camps held: name, post, dates, hours, role;
\item totals of persons, hours and camps.
\end{itemize}

It is signed by the candidate. A candidate from an institute also needs the signature of the institute's authorised signatory. An open candidate needs no institute signature.

\section*{Exemption}

The guideline gives an exemption in this component only to holders of a postgraduate yoga degree from an accredited institute who sit the Yoga Master certification. It does not apply at this level.

\section*{Start early}

Start collecting the record during the course.
\begin{itemize}
\item Teach a few real classes, for example at home, in a park, in a society or in a corporate setting.
\item Take a photograph.
\item Get a trainee's signed feedback form.
\item Keep a log of dates and hours.
\end{itemize}

The forms are the Board's. Use its current format, which is on its site and with the institute. The layout may change with a new version of the guideline.

\begin{recall}
\begin{itemize}
\item Marks: 10, part of the practical (140 in all).
\item Show: report with photographs, feedback form (Annexure A), experience format (Annexure B).
\item Modes of training: home based, park, society, corporate, other.
\item Open candidate: no institute signature.
\item No exemption at this level.
\end{itemize}
\end{recall}

\subsection*{Worked examples}

\textbf{What must a candidate show for field experience, and how many marks does it carry?} Field experience carries 10 marks, part of the practical (140 marks in all). The marking points are the report with photographs, the trainee's feedback form (Annexure A) and the experience details in the Board's format (Annexure B).

\textbf{A candidate builds the record during a 3 month part time course.}
\begin{itemize}
\item Teach a few real classes, for example at home, in a park, in a society or in a corporate setting.
\item Take a photograph at each programme for the report.
\item Ask a trainee to fill and sign the feedback form with the name and contact details, period and days of training, mode of training, reason for taking up yoga, feedback about the trainer and a mark out of 10.
\item Keep a log of dates and hours, so that the experience format can be filled with the persons trained, the mode, the totals of persons and hours, and the signatures.
\end{itemize}

\textbf{Who signs Annexure B.} A candidate from an institute signs, and the institute's authorised signatory also signs. An open candidate needs no institute signature. Use this as a quick check before submitting the form.

\subsection*{Common questions}

\begin{itemize}
\item \emph{When should I start collecting this record?} During the course.
\item \emph{Does the Board say exactly what the marks are given for?} The guideline gives the marks and the documents to show. The detailed content is not given. Ask the institution what it assesses.
\item \emph{Is a class log compulsory?} The Board asks for the report, the feedback form and the experience format. A short log of classes is a study suggestion and a practical way to keep the record.
\item \emph{Can I be exempt?} Not at this level.
\end{itemize}

\subsection*{Check yourself}

\begin{enumerate}
\item How many marks are given for field experience?
\item Fill in the blank: the feedback form is Annexure \underline{\hspace{1cm}} and the experience format is Annexure \underline{\hspace{1cm}}.
\item True or false: the postgraduate exemption in this component applies at Level 1.
\item Name four things the feedback form asks of the trainee.
\item List the modes of training named in the feedback form.
\item What does the experience format list besides the persons trained?
\item What must the candidate show for field experience?
\item A candidate has photographs and a signed feedback form but has not filled the experience details. Is the record complete?
\item A student asks where to get the forms and whether the layout may change. What do you say?
\end{enumerate}

\subsection*{Answers}

\begin{enumerate}
\item Ten, part of the practical (140 marks in all).
\item A for the feedback form; B for the experience format.
\item False. It is given only to holders of a postgraduate yoga degree from an accredited institute who sit the Yoga Master certification.
\item Any four of: name and contact details, period of training and number of days, mode of training, reason for taking up yoga, feedback about the trainer, a mark out of 10.
\item Home based, park, society, corporate, other.
\item Working experience with an organisation, yoga camps held, and totals of persons, hours and camps.
\item A report with photographs of a training programme conducted, a trainee's feedback form, and the experience details in the Board's format.
\item No. Annexure B is still needed.
\item The forms are the Board's. Use its current format, which is on its site and with the institute. The layout may change with a new version.
\end{enumerate}

\paragraph{Sources.} Yoga Certification Board, \emph{Guideline for Yoga Professionals}, the clause on the practical and Annexures A and B; Board syllabus for Level 1, Yoga Protocol Instructor; the official Ministry of AYUSH and Quality Council of India guidebook (older scheme, coverage check). In own words.
